\documentclass{article}
\usepackage[T1]{fontenc}
\usepackage{lmodern}
\usepackage{graphicx}
\usepackage{amsmath,amssymb}
\usepackage[a4paper,margin=2cm]{geometry}
\usepackage[absolute,overlay]{textpos}
\setlength{\TPHorizModule}{1mm}
\setlength{\TPVertModule}{1mm}
\usepackage[indonesian]{babel}
\usepackage{hyperref}
\setlength{\emergencystretch}{2em}
% unit: Halaman 49

\begin{document}

% === Halaman 49 (python/native) ===

Keamanan RC4

• RC4 memiliki kelemahan, yaitu byte-byte pada keystream awal pembangkitan memiliki 

korelasi yang tinggi dengan beberapa byte awal kunci

• Oleh karena itu direkomendasikan membuang 256 sampai 512 byte keystream awal


• RC4 adalah cipher alir, maka ia tidak kuat terhadap serangan seperti flip-bit attack maupun 

serangan-serangan stream attack lainnya. 

• Saat ini keamanan RC4 sudah berhasil dipecahkan dalam hitungan jam atau hari. 

• Pada bulan Februari 2015, RC4 dilarang penggunaannya di dalam Transport Layer Security 

(TLS) seperti disebutkan di dalam RFC 7465 (https://tools.ietf.org/html/rfc7465).

• Beberapa varian RC4 telah dibuat untuk mengatasi kelemahannya, yaitu Spritz, RC4A, 

VMPC, dan RC4+. 

49

\end{document}
